\documentclass[11pt]{article}
\usepackage[margin=1in]{geometry}
\usepackage{amsmath,amssymb,amsthm,booktabs,array}
\usepackage[T1]{fontenc}
\usepackage{lmodern}
\usepackage[hidelinks]{hyperref}
\newtheorem{theorem}{Theorem}
\newtheorem{lemma}[theorem]{Lemma}
\newtheorem{corollary}[theorem]{Corollary}
\title{Hardin's six-column trace recurrence\\and Barker's generating functions}
\author{Haobo Ma\\\small ChronoAI Pte Ltd, Singapore
\and Wenlin Zhang\\\small National University of Singapore\\\small ChronoAI Pte Ltd, Singapore}
\date{November 14, 2026}
\begin{document}
\maketitle
\begin{abstract}
We prove the empirical recurrence in R.\ H.\ Hardin's OEIS entry
A185765, revision 9, sole FORMULA field (October 3, 2025; submitted
February 3, 2011). It counts labeled $(n+1)\times6$ binary arrays whose
contiguous $2\times2$ blocks have unequal traces at horizontal and vertical
neighbors. A row-pair transfer operator reduces the all-length assertion to
three exact finite identities. The recurrence holds for every $n\geq13$
and fails at $n=12$, where its residual is $-4160$. Its reduced rational
generating function proves that nine is the least possible eventual
recurrence order over $\mathbb Q$. The same method proves the smaller-width
recurrences and Colin Barker's conjectural generating functions in
A185761--A185764, including his two-column closed form. Transposition,
divisibility by eight, and rationality at every fixed width give further
consequences for Hardin's bivariate table A185769.
\end{abstract}

\section{The counting question}
Hardin's entry~\cite{hardin765} records an empirical order-nine recurrence
for trace-avoiding arrays at every index beyond twelve. We establish
that recurrence with its printed threshold and its least eventual order.
The related entries~\cite{hardin761,hardin762,hardin763,hardin764} record
smaller-width empirical recurrences and Barker's conjectural generating
functions, added in February and April 2018. We prove these formulas by
the same finite-state argument. Their attribution remains Hardin's for
the counting sequences and empirical recurrences, and Barker's for the
generating-function and closed-form conjectures.
The finite-state transfer method is classical; see Stanley~\cite[Section~4.7]{stanley}.
Here the specified recurrence polynomials follow from explicit finite
identities and linear propagation.
For $n\geq1$, let $a_n$ count matrices
$A\in\{0,1\}^{(n+1)\times6}$ satisfying the following conditions.
With rows numbered $0,\ldots,n$ and columns $0,\ldots,5$, set
\[
 t_{i,j}=A_{i,j}+A_{i+1,j+1}\quad(0\leq i<n,\ 0\leq j<5).
\]
Require $t_{i,j}\ne t_{i,j+1}$ whenever $j<4$, and
$t_{i,j}\ne t_{i+1,j}$ whenever $i<n-1$.
Thus the trace is the ordinary integer diagonal sum, not the sum of all
four entries or addition modulo two. Boundaries are open, and arrays are
not identified under symmetries.

\begin{theorem}\label{thm:main}
For every integer $n\geq13$,
\begin{align}
 a_n={}&a_{n-1}+30a_{n-2}+24a_{n-3}-95a_{n-4}-15a_{n-5}\notag\\
       &+64a_{n-6}-116a_{n-7}+16a_{n-8}+16a_{n-9}.\label{eq:rec}
\end{align}
The difference between the left and right sides at $n=12$ is $-4160$.
\end{theorem}

\section{Arrays as walks}
Write $R=\{0,1\}^6$ and, for $r,s\in R$, put
$\tau_j(r,s)=r_j+s_{j+1}$ for $0\leq j<5$.
The state set is
\[
 S=\{(r,s)\in R^2:\tau_j(r,s)\ne\tau_{j+1}(r,s)
                         \text{ for }0\leq j<4\}.
\]
An edge $(r,s)\longrightarrow(s,t)$ exists precisely when $(s,t)\in S$
and $\tau_j(r,s)\ne\tau_j(s,t)$ for every $0\leq j<5$.
For a vector $v\in\mathbb Z^S$ define
\[
 (Tv)_{(r,s)}=\sum_{(r,s)\to(s,t)}v_{(s,t)},
 \qquad \mathbf1=(1)_{s\in S}.
\]
There are $680$ states and $3456$ edges; these numbers are obtained by
enumerating the $64^2$ row pairs and $64^3$ row triples.

\begin{lemma}\label{lem:count}
For every $n\geq1$, $a_n=\mathbf1^{\mathsf T}T^{n-1}\mathbf1$.
\end{lemma}
\begin{proof}
An array with rows $r_0,\ldots,r_n$ determines the walk
$(r_0,r_1),\ldots,(r_{n-1},r_n)$. Membership of each state in $S$
is exactly the horizontal trace condition. Each successive edge is
exactly the vertical trace condition between the two corresponding
rows of blocks. Conversely, overlapping row pairs in such a walk
reconstruct the unique array, so this correspondence is a bijection.
The coordinate $(T^k\mathbf1)_{(r,s)}$ counts walks of $k$ edges
starting at $(r,s)$: it is one at $k=0$, and the defining sum for $T$
is the induction step. Summing over the first state with $k=n-1$
proves the formula.
\end{proof}

\section{The finite identities and their propagation}
Define
\begin{align*}
 Q(X)={}&X^9-X^8-30X^7-24X^6+95X^5+15X^4\\
        &-64X^3+116X^2-16X-16,\qquad
 w_k=T^kQ(T)\mathbf1.
\end{align*}
The finite identities needed below are
\begin{equation}\label{eq:finite}
 w_5=0,\qquad \mathbf1^{\mathsf T}w_3=0,\qquad
 \mathbf1^{\mathsf T}w_4=0.
\end{equation}
Table~\ref{tab:finite} gives the full check's scalar summaries. In
particular, the zero nonzero-coordinate count in its final row asserts
the vector identity, not just cancellation of its coordinate sum.

\begin{table}[ht]
\centering
\begin{tabular}{rrrr}
\toprule
$k$ & Nonzero coordinates of $w_k$ & $\sum_{s\in S}(w_k)_s$
    & $\max_{s\in S}|(w_k)_s|$\\
\midrule
0 & 680 & $-23144$ & 9540\\
1 & 680 & 27200 & 23192\\
2 & 680 & $-4160$ & 14932\\
3 & 680 & 0 & 10640\\
4 & 376 & 0 & 344\\
5 & 0 & 0 & 0\\
\bottomrule
\end{tabular}
\caption{Exact integer evaluation of the six residual vectors.}\label{tab:finite}
\end{table}

Here is a complete finite algorithm for~\eqref{eq:finite}. Order binary
rows lexicographically and retain exactly the pairs in $S$. Construct
each outgoing edge list by testing every $t\in R$. Starting with
$v_0=\mathbf1$, calculate $v_{j+1}=Tv_j$ through $j=13$ and form
\[
 (w_k)_s=(v_{k+9})_s-\sum_{d=1}^9 c_d(v_{k+9-d})_s
 \quad(0\leq k\leq5),
\]
where $(c_1,\ldots,c_9)=(1,30,24,-95,-15,64,-116,16,16)$.
The supplied \texttt{verify.py} implements precisely these bounded
loops with integer arithmetic. It also evaluates $Q(T)\mathbf1$ by
Horner's rule and compares all coordinates with the expanded expression.
An independent bit-row encoding checks every state and outgoing edge.
Literal enumeration of all two-row and three-row arrays gives $680$
and $3456$, respectively, and checks the indexing against the array
definition. Thus~\eqref{eq:finite} is an exhaustive finite computation;
the passage to all lengths is the following written argument.

\begin{proof}[Proof of Theorem~\ref{thm:main}]
Linearity gives $w_{k+1}=Tw_k$ for all $k\geq0$.
Consequently $w_5=0$ implies $w_k=0$ for every $k\geq5$ by induction.
Together with the two scalar equalities in~\eqref{eq:finite}, this gives
$\mathbf1^{\mathsf T}w_k=0$ for all $k\geq3$.
For $n\geq13$, put $k=n-10\geq3$. Lemma~\ref{lem:count} then gives
\[
 0=\mathbf1^{\mathsf T}w_k
   =a_{k+10}-\sum_{d=1}^9c_da_{k+10-d},
\]
which is~\eqref{eq:rec}. At $n=12$ the same identity for the residual
uses $k=2$, whose sum in Table~\ref{tab:finite} is $-4160$.
\end{proof}

The first twelve values, sufficient to start the recurrence at $n=13$, are
\begin{gather*}
680,\ 3456,\ 19992,\ 118248,\ 725352,\ 4430912,\\
27172072,\ 166155112,\ 1017804600,\ 6228062600,\\
38135220896,\ 233413737792.
\end{gather*}
Forward walk counts agree with all $200$ terms of Hardin's linked
table~\cite{hardin765}. This agreement is a finite comparison; the
all-length conclusion follows from~\eqref{eq:finite} and the proof above.

\section{A rational generating function and least eventual order}
\begin{corollary}\label{cor:minimal}
Let $F_6(z)=\sum_{n\geq1}a_nz^n$ and
$D_6(z)=1-z-30z^2-24z^3+95z^4+15z^5-64z^6+116z^7-16z^8-16z^9$.
Then
\begin{equation}\label{eq:gf6}
 F_6(z)=\frac{8zN_6(z)}{D_6(z)},
\end{equation}
where
\begin{align*}
N_6(z)={}&85+347z-483z^2-2718z^3-1375z^4+2104z^5\\
 &+6276z^6+4796z^7-7734z^8-2893z^9+3400z^{10}-520z^{11}.
\end{align*}
No homogeneous constant-coefficient recurrence over $\mathbb Q$ of
order less than nine holds for all sufficiently large $n$.
\end{corollary}
\begin{proof}
Multiplying the first twelve terms by $D_6$ gives $8zN_6$;
all coefficients from degree thirteen onward vanish by
Theorem~\ref{thm:main}. This proves~\eqref{eq:gf6} as a formal power-series
identity. The Euclidean algorithm applied to $N_6,D_6$ modulo $3$ gives
successive nonzero remainder degrees $8,6,5,4,3,1,0$, ending in the
constant $2$. Explicitly, their coefficient lists in increasing degree are
\[
\begin{gathered}
(2,0,0,1,1,2,0,2,1),\quad(2,1,0,2,2,2,1),\\
(0,2,1,1,2,1),\quad(2,1,1,1,1),\quad(1,2,2,2),\quad(2,2),\quad(2).
\end{gathered}
\]
Both leading coefficients are nonzero modulo $3$. By Gauss's lemma a
nonconstant common factor over $\mathbb Q$ would therefore remain a
nonconstant common factor modulo $3$, a contradiction. Also $D_6(0)=1$,
so $8zN_6$ and $D_6$ are coprime over $\mathbb Q$.
An eventual recurrence with denominator polynomial $E(z)$ would make
$E(z)F_6(z)$ a polynomial. Coprimality forces $D_6\mid E$, hence
$\deg E\geq9$. Theorem~\ref{thm:main} attains this bound.
\end{proof}

\section{Smaller widths and Barker's conjectures}
For $m\geq2$ let $a^{(m)}_n$ count the same arrays with $m$ columns;
in particular $a^{(6)}_n=a_n$. Define $S_m,T_m$ exactly as in Section~2,
using binary rows of length $m$ and $m-1$ traces. The same bijection gives
$a^{(m)}_n=\mathbf1^{\mathsf T}T_m^{n-1}\mathbf1$.

\begin{corollary}\label{cor:siblings}
For each row of Table~\ref{tab:siblings},
\[
 a^{(m)}_n=\sum_{j=1}^{d_m}c_j a^{(m)}_{n-j}\quad(n\geq s_m),
 \qquad
 \sum_{n\geq1}a^{(m)}_nz^n=\frac{8zN_m(z)}{1-\sum_{j=1}^{d_m}c_jz^j}.
\]
The numerators are
\begin{align*}
 N_2&=2+3z, &N_3&=5+10z+z^2,\\
 N_4&=13+33z-16z^3+7z^4-z^5,\\
 N_5&=33+108z-43z^2-249z^3-69z^4+5z^5+43z^6+81z^7+18z^8.
\end{align*}
These are Barker's formulas in A185761--A185764~\cite{hardin761,hardin762,hardin763,hardin764}.
\end{corollary}
\begin{table}[ht]
\centering\small
\begin{tabular}{rrrrlrr}
\toprule
$m$ & $|S_m|$ & Edges & $d_m$ & $(c_1,\ldots,c_{d_m})$ & $s_m$ & $h_m$\\
\midrule
2 & 16 & 40 & 2 & $(1,4)$ & 3 & 1\\
3 & 40 & 120 & 3 & $(1,6,1)$ & 4 & 2\\
4 & 104 & 368 & 5 & $(1,10,3,-5,1)$ & 7 & 3\\
5 & 264 & 1128 & 7 & $(1,17,9,-22,-1,1,-6)$ & 10 & 3\\
\bottomrule
\end{tabular}
\caption{Smaller-width recurrences and vector extinction exponents.}\label{tab:siblings}
\end{table}
\begin{proof}
Put $Q_m(X)=X^{d_m}-\sum_jc_jX^{d_m-j}$ and
$w^{(m)}_k=T_m^kQ_m(T_m)\mathbf1$.
The same exhaustive algorithm as in Section~3 gives
$w^{(m)}_{h_m}=0$. The additional scalar identities
$\mathbf1^{\mathsf T}w^{(m)}_k=0$ are needed at
$k=0$ for $m=2$; $k=0,1$ for $m=3$; $k=1,2$ for $m=4$; and $k=2$
for $m=5$. These are exact finite identities in the displayed state spaces,
reproduced by \texttt{verify.py}. For $k\geq h_m$, vector vanishing
propagates under $T_m$. Setting $n=k+d_m+1$ proves each recurrence
at every $n\geq s_m$. Multiplication by its displayed denominator then
leaves only degrees below $s_m$; exact counts at these initial indices
give $8zN_m$ and prove each generating function.
\end{proof}

\begin{corollary}\label{cor:closed}
For $n\geq1$, with $r_\pm=(1\pm\sqrt{17})/2$,
\[
 a^{(2)}_n=
 \frac{(-13+3\sqrt{17})r_-^n+(13+3\sqrt{17})r_+^n}{\sqrt{17}}.
\]
\end{corollary}
\begin{proof}
Both $r_\pm$ satisfy $r^2=r+4$, so the displayed expression satisfies
the two-column recurrence. Its values at $n=1,2$ are $16,40$, respectively,
the counts from Table~\ref{tab:siblings}. Induction proves equality
with $a^{(2)}_n$ for all $n\geq1$. This is Barker's closed form in
the FORMULA field of A185761, February 22, 2018~\cite{hardin761}.
\end{proof}

\section{The cross-referenced bivariate table}
Let $b_{n,k}$ count admissible $(n+1)\times(k+1)$ arrays under the same
trace rule; this is Hardin's A185769~\cite{hardin769}.
\begin{corollary}\label{cor:transpose}
For all $n,k\geq1$, $b_{n,k}=b_{k,n}$. In particular,
$b_{n,5}=b_{5,n}=a_n$, so both the fifth column and fifth row of
A185769 satisfy~\eqref{eq:rec}, with the varying index at least $13$.
\end{corollary}
\begin{proof}
Transposition sends a block at $(i,j)$ to the block at $(j,i)$ and
preserves its diagonal sum. It exchanges horizontal and vertical
adjacency, so it is a bijection between the two sets of admissible
arrays. The six-column specialization is $k=5$, giving the result.
\end{proof}

\begin{corollary}\label{cor:general}
For all $n,k\geq1$, $8\mid b_{n,k}$. For each fixed $k\geq1$,
the generating function of $b_{n,k}$ in $n$ is rational:
\[
 \sum_{n\geq1}b_{n,k}z^n
 =z\mathbf1^{\mathsf T}(I-zT_{k+1})^{-1}\mathbf1.
\]
Consequently every fixed column, and by transposition every fixed row,
satisfies some eventual homogeneous constant-coefficient recurrence.
\end{corollary}
\begin{proof}
The top-right and bottom-left entries of the array occur in no block
trace, giving four free choices. Complementing all remaining entries
sends every block trace $t$ to $2-t$, so it pairs admissible assignments
of the remaining entries without fixed points. There is at least one
remaining entry for $n,k\geq1$; hence the total is divisible by eight.
For rationality, sum the array--walk identity over $n\geq1$ and use
$\sum_{j\geq0}z^jT_{k+1}^j=(I-zT_{k+1})^{-1}$ as a formal matrix
power series. The adjugate formula makes each entry rational, since
$\det(I-zT_{k+1})$ has constant term $1$. Multiplying by this determinant
gives an eventual recurrence. Transposition proves the row statement.
\end{proof}

\section{Source coverage and verification}
\noindent
\begin{tabular}{>{\raggedright\arraybackslash}p{0.34\linewidth}>{\raggedright\arraybackslash}p{0.58\linewidth}}
\toprule
Source item & Coverage\\
\midrule
A185765, revision 9, sole FORMULA & Proved for every $n>12$ by Theorem~\ref{thm:main}.\\
A185765, COMMENT & Its identification with column five of A185769
follows from the definitions.\\
A185765, four examples and term table & Four examples satisfy the rule;
all $200$ tabulated terms agree with exact walk counts.\\
A185761, revision 8, FORMULA & Hardin's recurrence and Barker's generating
function and closed form proved by Corollaries~\ref{cor:siblings}--\ref{cor:closed}.\\
A185762--A185764, revisions 9, 9, 10, FORMULA & Hardin's recurrences and
Barker's generating functions proved by Corollary~\ref{cor:siblings}.\\
A185769, revision 9, COMMENTS and examples & All 70 displayed counts
and seven examples checked; symmetry proved by Corollary~\ref{cor:transpose}.\\
Least eventual order nine; divisibility; fixed-width rationality & Not posed
in these entries; proved by Corollaries~\ref{cor:minimal} and~\ref{cor:general}.\\
Extension to $n=12$ & Not asserted in the source; fails with residual $-4160$.\\
Larger-width specific formulas, e.g.\ \href{https://oeis.org/A185766}{A185766}
and \href{https://oeis.org/A185768}{A185768} & Not covered. General rationality
does not establish their particular coefficients or thresholds.\\
\bottomrule
\end{tabular}


\paragraph{Verification and availability.}
The array--walk bijection, propagation, minimal-order argument and
corollaries have written proofs above. Python 3's standard-library
\texttt{verify.py} reproduces the finite vector and scalar identities,
generating-function coefficients and polynomial remainder lists.
It also compares all $200$ A185765 terms with the linked data when
that file is supplied, and checks all $70$ displayed A185769 counts.

For six columns the pinned Lean file
\href{https://github.com/the-omega-institute/trureturing/blob/66c49bfe9a1295ceea872b6e6a1d8146fb351237/D5/S3/Combinatorics/TraceAvoidingArrays/HardinA185765Transfer.lean}{\texttt{HardinA185765Transfer.lean}}
contains \texttt{counting\_bridge}: the literal natural-number array
cardinality at $n=k+1$, embedded in $\mathbb Z$, equals the sum of its
continuation counts $u_k$ over all $64^2$ binary row pairs. Pairs outside
$S$ have zero continuation count. The file
\href{https://github.com/the-omega-institute/trureturing/blob/66c49bfe9a1295ceea872b6e6a1d8146fb351237/D5/S3/Combinatorics/TraceAvoidingArrays/HardinA185765Recurrence.lean}{\texttt{HardinA185765Recurrence.lean}}
contains \texttt{finite\_certificate}, establishing the vector identity
at exponent five and scalar identities at exponents three and four,
and \texttt{result}, giving~\eqref{eq:rec} for every $n\geq13$.
Both namespaces start with
\texttt{D5.S3.Combinatorics.}\allowbreak\texttt{TraceAvoidingArrays}.
These declarations cover the six-column counting bridge and recurrence;
the generating functions, minimal order, smaller-width corollaries,
transposition and divisibility are the written arguments and finite
identities specified above.

\clearpage
\begin{thebibliography}{7}
\bibitem{hardin761}
R.\ H.\ Hardin, \emph{Number of $(n+1)\times2$ binary arrays with no
$2\times2$ subblock trace equal to any horizontal or vertical neighbor
$2\times2$ subblock trace}, OEIS A185761, February 3, 2011;
revision 8, February 22, 2018. FORMULA: generating-function and
closed-form conjectures by Colin Barker, February 22, 2018.
\url{https://oeis.org/A185761}.
\bibitem{hardin762}
R.\ H.\ Hardin, \emph{Number of $(n+1)\times3$ binary arrays with no
$2\times2$ subblock trace equal to any horizontal or vertical neighbor
$2\times2$ subblock trace}, OEIS A185762, February 3, 2011;
revision 9, April 16, 2018. FORMULA: empirical generating function
by Colin Barker, April 16, 2018. \url{https://oeis.org/A185762}.
\bibitem{hardin763}
R.\ H.\ Hardin, \emph{Number of $(n+1)\times4$ binary arrays with no
$2\times2$ subblock trace equal to any horizontal or vertical neighbor
$2\times2$ subblock trace}, OEIS A185763, February 3, 2011;
revision 9, April 16, 2018. FORMULA: empirical generating function
by Colin Barker, April 16, 2018. \url{https://oeis.org/A185763}.
\bibitem{hardin764}
R.\ H.\ Hardin, \emph{Number of $(n+1)\times5$ binary arrays with no
$2\times2$ subblock trace equal to any horizontal or vertical neighbor
$2\times2$ subblock trace}, OEIS A185764, February 3, 2011;
revision 10, April 16, 2018. FORMULA: empirical generating function
by Colin Barker, April 16, 2018. \url{https://oeis.org/A185764}.
\bibitem{hardin765}
R.\ H.\ Hardin, \emph{Number of $(n+1)\times6$ binary arrays with no
$2\times2$ subblock trace equal to any horizontal or vertical neighbor
$2\times2$ subblock trace}, OEIS A185765, submitted February 3, 2011,
revision 9, October 3, 2025. \url{https://oeis.org/A185765}.
Term table: \url{https://oeis.org/A185765/b185765.txt}.
\bibitem{hardin769}
R.\ H.\ Hardin, \emph{$T(n,k)$: number of $(n+1)\times(k+1)$ binary
arrays with no $2\times2$ subblock trace equal to any horizontal or vertical
neighbor $2\times2$ subblock trace},
OEIS A185769, submitted February 3, 2011, revision 9, October 3, 2025.
\url{https://oeis.org/A185769}.
\bibitem{stanley}
R.\ P.\ Stanley, \emph{Enumerative Combinatorics}, Volume~1, second edition,
Cambridge University Press, 2011, Section~4.7, ``The Transfer-Matrix Method.''
Author's book page: \url{https://math.mit.edu/~rstan/ec/ec1/}.
\end{thebibliography}
\smallskip
\noindent\small \copyright\ 2026 Haobo Ma and Wenlin Zhang.
Licensed under Creative Commons Attribution 4.0 International (CC BY 4.0).
\end{document}
